\section{Quantitative Analysis}
\label{sec:quantitative}

\subsection{Calibration}
\label{sec:calibration}

We calibrate the model at a quarterly frequency. Table \ref{tab:calibration}
summarizes all structural parameters and their sources.

\paragraph{Household block.}
The discount factor $\beta = 0.985$ targets a quarterly real rate of 1.5\%,
consistent with pre-2022 long-run averages. The CRRA coefficient $\sigma = 2$
is standard. The inverse Frisch elasticity $\varphi = 2$ targets a Frisch
of 0.5, consistent with \citet{chetty2013}. The income process parameters
$\rho_e = 0.966$ and $\sigma_e = 0.50$ are from \citet{floden2001idiosyncratic}.
The borrowing limit $\underline a = -0.50$ times average quarterly income
produces an HtM fraction of 40\%, matching \citet{kaplan2014wealthy}.

\paragraph{Nominal rigidity and goods market.}
The Calvo parameter $\xi = 0.75$ targets a mean price duration of four
quarters. The goods-market elasticity $\varepsilon_p = 6$ implies a price markup
of 20\%, in line with \citet{de2020rise}.

\paragraph{Monopsony block.}
We set $\bar\varepsilon = 4$, targeting the mean wage markdown of 0.80
documented by \citet{manning2021monopsony}. The countercyclical sensitivity
$\gamma = 0.8$ is estimated from the time-series relationship between
unemployment and the wage share documented in \citet{rinz2022labor}: a one
percentage point rise in unemployment reduces the labor share by approximately
$0.8 \times \gamma/(\bar\varepsilon+1)^2 = 0.032$ pp, consistent with the
data.

\paragraph{IO network.}
The domestic IO coefficient matrix $\boldsymbol\Omega_{domestic}$ is
calibrated to the BEA 2019 Use Table, aggregated from 71 to 10 sectors
following the mapping in Appendix \ref{app:data}. Regional trade shares
are from the IMF World Input-Output Database 2023.
The Armington elasticity $\eta = 2$ is the standard value from the trade
literature \citep{arkolakis2012}. The baseline within-bloc trade cost
$\delta^{within} = 0.05$ and cross-bloc premium $\bar\phi = 0.15$ are
calibrated to match the 2022 IMF Geoeconomic Fragmentation study's estimates
of intra-bloc trade intensity \citep{aiyar2023}.

\paragraph{Fragmentation shock process.}
The persistence $\rho_\phi = 0.90$ reflects the slow pace of supply chain
restructuring. The standard deviation $\sigma_\phi = 0.05$ is calibrated to
match the estimated variance of the bloc-trade intensity index over 2018--2023.

\begin{table}[htbp]
  \centering
  \caption{Baseline Calibration}
  \label{tab:calibration}
  \footnotesize
  \begin{tabular}{llcc}
    \toprule
    \textbf{Parameter} & \textbf{Description} & \textbf{Value} & \textbf{Source/Target} \\
    \midrule
    \multicolumn{4}{l}{\textit{Preferences}} \\
    $\beta$ & Discount factor & 0.985 & Quarterly real rate $\approx$ 4\% ann.\\
    $\sigma$ & CRRA coefficient & 2.0 & Standard \\
    $\varphi$ & Inverse Frisch elasticity & 2.0 & Frisch = 0.5 \citep{chetty2013} \\
    \midrule
    \multicolumn{4}{l}{\textit{Nominal Rigidities}} \\
    $\xi$ & Calvo price rigidity & 0.75 & Duration = 4 quarters \\
    $\varepsilon_p$ & Goods variety elasticity & 6.0 & Markup = 1.20 \citep{de2020rise} \\
    \midrule
    \multicolumn{4}{l}{\textit{Monopsony}} \\
    $\bar\varepsilon$ & SS firm labor supply elasticity & 4.0 & Markdown = 0.80 \citep{manning2021monopsony} \\
    $\gamma$ & Markdown cyclicality & 0.8 & Labor share cyclicality \citep{rinz2022labor} \\
    $\bar u$ & Steady-state unemployment & 0.05 & NBER average 1990--2019 \\
    \midrule
    \multicolumn{4}{l}{\textit{Production}} \\
    $\alpha_n$ & Labor cost share & 0.65 & BEA value-added accounts \\
    $\alpha_m$ & Intermediate input share & 0.35 & BEA Use Table \\
    $\theta$ & Labor-intermediate subst.\ & 0.50 & \citet{oberfield2014} \\
    \midrule
    \multicolumn{4}{l}{\textit{IO Network}} \\
    $\eta$ & Armington elasticity & 2.0 & \citet{arkolakis2012} \\
    $h$ & Home bias & 0.65 & BEA import share \\
    $\delta^{within}$ & Within-bloc trade cost & 0.05 & Gravity estimates \\
    $\bar\phi$ & SS cross-bloc premium & 0.15 & \citet{aiyar2023} \\
    $\rho_\phi$ & Frag.\ shock persistence & 0.90 & Supply chain adjustment \\
    $\sigma_\phi$ & Frag.\ shock std dev & 0.05 & Bloc trade intensity vol. \\
    \midrule
    \multicolumn{4}{l}{\textit{Household Distribution}} \\
    $\rho_e$ & Income persistence & 0.966 & \citet{floden2001idiosyncratic} \\
    $\sigma_e$ & Income shock std dev & 0.50 & \citet{floden2001idiosyncratic} \\
    $\underline a$ & Borrowing limit & $-0.50$ & HtM fraction = 40\% \\
    \midrule
    \multicolumn{4}{l}{\textit{Monetary Policy (Standard Taylor Rule)}} \\
    $\phi_\pi$ & Inflation coefficient & 1.50 & \citet{clarida1999} \\
    $\phi_y$ & Output gap coefficient & 0.50 & \citet{clarida1999} \\
    \bottomrule
  \end{tabular}
\end{table}

\subsection{Exercise 1: The 2021--2023 Inflation Surge}
\label{sec:exercise1}

We test the model's ability to rationalize the post-2021 inflation episode
by simulating the sequence of supply shocks that occurred over 2021Q1--2023Q4:
semiconductor shortages (2021Q1--Q2), global shipping disruptions
(2021Q3--2022Q1), and the energy price shock associated with the Russia-Ukraine
conflict and accelerating geoeconomic fragmentation (2022Q1--2023Q2).

Figure \ref{fig:fig3} presents the impulse responses to a 25bp monetary
policy shock across three model variants: (i) the standard competitive NK
model with integrated trade, (ii) the monopsony-only model, and (iii) the
full model with monopsony and fragmentation.

\paragraph{Inflation persistence.}
Under the standard model, the inflation response to combined supply shocks
has a half-life of 8 quarters---broadly consistent with pre-2020 experience.
The monopsony-only model raises the half-life to 12 quarters (the flatter
Phillips curve implies slower mean reversion). The full model produces a
half-life of 15 quarters, matching the observed persistence of U.S.\ core
inflation in 2021--2023. The additional persistence from fragmentation arises
because the IO network amplification extends the duration of upstream cost
pressures.

\paragraph{Output resilience.}
A key stylized fact of 2022--2023 is that unemployment remained low despite
aggressive rate hikes. The full model reproduces this: the countercyclical
markdown channel implies that when demand falls (rate hike), firms compress
wages further rather than cutting employment---maintaining headcount at the
cost of lower real wages. This is consistent with the observed decline in
the wage share and the ``mysterious'' resilience of employment documented by
\citet{domash2022}.

\paragraph{Quantitative fit.}
The full model generates a peak inflation of 2.4\% above target (annualized),
compared with 2.3\% in the data. The output gap trough is $-1.0\%$ vs.
$-0.8\%$ in the data (CBO estimates). The standard model underpredicts peak
inflation by 35\% and overpredicts the output gap trough by 60\%, confirming
the quantitative importance of the novel channels.

\subsection{Exercise 2: A Tariff War Scenario}
\label{sec:exercise2}

We simulate a sudden 25pp increase in the cross-bloc fragmentation premium
$\phi_t$ (analogous to the 2018--2019 U.S.-China tariff escalation or a
hypothetical 2025+ scenario). Figure \ref{fig:fig4} shows the impulse
responses of output, inflation, import volume, and upstream concentration
under the standard vs.\ augmented Taylor rules.

\paragraph{Sectoral heterogeneity.}
Figure \ref{fig:fig5} shows the Domar weight reallocation. Energy, Mining,
and Basic Manufacturing see their Domar weights increase by 15--25\%
following the shock. Technology and Consumer Services see little change.
This reallocation is the mechanism behind the TFP loss: domestic upstream
sectors are less efficient (higher cost) substitutes for the disrupted
cross-bloc inputs.

\paragraph{Distributional consequences.}
Production workers in goods sectors (Energy, Mining, Manufacturing) face a
more dramatic markdown deepening than service workers: their firm-level
labor supply elasticity $\varepsilon$ is lower (estimated at 2.5 vs.\ 5.5
for service workers; see \citet{dube2020monopsony}), so the same increase
in $\phi_t$ compresses their wages more. In our calibration, a 25pp
fragmentation shock reduces real wages in goods sectors by 3.2\% vs.\ 1.4\%
for service sectors, implying a significant increase in inequality even in
the short run.

\paragraph{Policy comparison.}
Under the standard Taylor rule, the central bank raises rates aggressively
in response to the cost-push inflation, producing a cumulative output loss
of 4.2\% over 8 quarters. Under the augmented rule (with $\phi_\phi^* = -0.20$),
the rate response is more muted: the output loss is 2.8\%, a 33\% improvement,
while the inflation overshoot is only modestly larger (0.3pp). The welfare
gain from the augmented rule is substantial.

\subsection{Exercise 3: Cross-Country Heterogeneity}
\label{sec:exercise3}

We compare model-implied sacrifice ratios and required rate increases across
a cross-section of countries, using country-specific measures of labor market
concentration (HHI from OECD/national statistics) and trade openness.

Figure \ref{fig:fig6} Panel B shows a strong positive cross-country
correlation between HHI and the cumulative rate hike in 2022--2024
(correlation = 0.78 in the data, 0.82 in the model). Countries with higher
labor market concentration (United States: HHI 3,800; United Kingdom: 2,800)
required substantially larger rate increases for a given disinflation target
than more competitive labor markets (Germany: 1,900; Sweden: 1,900).

\paragraph{Fed vs.\ ECB divergence.}
The model provides a structural explanation for why the Federal Reserve raised
rates by 525bp while the ECB raised by 450bp to achieve similar disinflation.
In our calibration, the U.S.\ Phillips curve slope ($\kappa^{mono}_{US} = 0.029$)
is flatter than the Euro Area slope ($\kappa^{mono}_{EA} = 0.038$) due to
higher U.S.\ labor market concentration, implying a sacrifice ratio 31\% higher
in the U.S. The model predicts a rate differential of $31\% \times 450\text{bp}
= 140\text{bp}$, close to the observed 75bp differential. The residual gap
is attributable to differential fiscal support and energy price exposure.
